\chapter{Coronaviruses}

\section*{Definition and Overview}
Coronaviruses are a large family of viruses that cause illness ranging from the common cold to more severe disease. They are named for the crown-like spikes on their surface. Seven coronaviruses are known to infect humans: four cause mild upper respiratory infections, while three others -- SARS-CoV, MERS-CoV, and SARS-CoV-2, the cause of COVID-19 -- can produce severe respiratory illness. Coronaviruses spread mainly through respiratory droplets and aerosols when an infected person coughs, sneezes, or talks. They have caused outbreaks and pandemics, most notably COVID-19, which began in 2019 and spread worldwide.

\section*{Epidemiology}
Coronaviruses are found worldwide and circulate continuously. The four common human coronaviruses cause a significant share of common colds, mostly in winter, and re-infections occur throughout life. SARS-CoV-2, the virus responsible for COVID-19, spread across the globe from early 2020, causing millions of deaths and widespread disruption before vaccines and natural immunity substantially reduced severe illness. MERS-CoV continues to cause sporadic cases in the Middle East, and SARS-CoV caused a contained outbreak in 2002--2003. New coronaviruses are expected to emerge from animal reservoirs, making surveillance important.

\section*{Aetiology and Pathophysiology}
Coronaviruses are enveloped RNA viruses that infect the cells lining the respiratory tract.

\begin{itemize}
    \item \textbf{Structure:} a crown of spike proteins that bind to receptors on host cells, such as ACE2 for SARS-CoV-2
    \item \textbf{Transmission:} respiratory droplets and aerosols, plus contact with contaminated surfaces
    \item \textbf{Entry:} the virus enters cells of the nose, throat, and airways and replicates
    \item \textbf{Immune response:} the body's defences clear the virus but can cause inflammation in severe disease
    \item \textbf{Severe disease:} the virus can reach the lungs, causing pneumonia, and trigger an excessive inflammatory response
    \item \textbf{Common cold coronaviruses:} cause limited, mild upper airway infection
    \item \textbf{High-impact coronaviruses:} SARS-CoV, MERS-CoV, and SARS-CoV-2 can cause severe pneumonia and systemic illness
\end{itemize}

\section*{Clinical Features}
Symptoms vary from none to life-threatening.

\begin{itemize}
    \item Runny nose, sneezing, sore throat, and cough (common cold pattern)
    \item Fever, chills, and fatigue
    \item Headache and muscle aches
    \item Loss of sense of smell or taste (common with COVID-19)
    \item Shortness of breath and chest tightness in more severe infection
    \item Gastrointestinal symptoms such as nausea, vomiting, or diarrhoea in some cases
    \item Severe complications include pneumonia, respiratory failure, and blood clotting problems
\end{itemize}

\section*{Differential Diagnosis}
Respiratory infections have many causes that are considered together.

\begin{itemize}
    \item Influenza and other respiratory viruses, such as RSV and adenovirus
    \item The common cold from rhinoviruses and other agents
    \item Bacterial pneumonia or bronchitis
    \item Allergic rhinitis or asthma
    \item Strep throat and other bacterial throat infections
    \item Other viral infections with systemic symptoms, such as glandular fever
\end{itemize}

\section*{When to Seek Medical Care}
Most coronavirus infections are mild and resolve with rest and fluids at home. Seek medical advice if symptoms are severe, persistent, or worsening, if you are at high risk (older age, chronic illness, pregnancy, or immune suppression), or if you need testing or treatment such as antiviral medicines.

\textbf{Call 000 (or local emergency services) for:}
\begin{itemize}
    \item Difficulty breathing, shortness of breath at rest, or chest pain
    \item Blue lips or face, or a severe drop in oxygen levels
    \item Confusion, drowsiness, or difficulty waking
    \item In children: difficulty breathing, poor feeding, or dehydration
    \item A sudden severe headache or neurological symptoms
\end{itemize}

\section*{Diagnosis and Investigations}
Most mild cases need no tests. Testing and investigations are used in specific situations.

\begin{itemize}
    \item \textbf{Rapid antigen test:} a nasal swab giving a result within minutes, used for COVID-19
    \item \textbf{PCR test:} a more sensitive nose and throat swab that detects viral RNA
    \item \textbf{Blood tests:} full blood count, inflammatory markers, and organ function in severe illness
    \item \textbf{Oxygen monitoring:} pulse oximetry to measure blood oxygen levels
    \item \textbf{Chest X-ray or CT:} in pneumonia, to assess lung involvement
    \item \textbf{Other respiratory testing:} to detect influenza and other viruses when relevant
\end{itemize}

\section*{Management}
Treatment depends on severity and risk.

\begin{itemize}
    \item \textbf{Self-care:} rest, fluids, and paracetamol or ibuprofen for fever and aches
    \item \textbf{Antivirals:} medicines for COVID-19 such as nirmatrelvir--ritonavir or molnupiravir, available to high-risk people
    \item \textbf{Influenza antivirals:} oseltamivir, where influenza is confirmed and treatment is early
    \item \textbf{Oxygen and hospital care:} for people with low oxygen levels or severe pneumonia
    \item \textbf{Anti-inflammatory treatment:} corticosteroids and other agents for severe inflammation
    \item \textbf{Blood-thinning medicines:} to prevent clotting complications in hospitalised patients
    \item \textbf{Infection control:} isolation, masks, and good hygiene to prevent spread
    \item \textbf{Vaccination:} COVID-19 and influenza vaccines reduce severe illness and hospitalisation
\end{itemize}

\section*{Complications}
\begin{itemize}
    \item Pneumonia and respiratory failure
    \item Acute respiratory distress syndrome (ARDS)
    \item Blood clots in the lungs, legs, or other organs
    \item Heart muscle inflammation (myocarditis) or abnormal heart rhythms
    \item Multi-organ failure in severe cases
    \item Long COVID: persistent fatigue, breathlessness, and cognitive symptoms after infection
    \item Secondary bacterial infections
    \item In children: multisystem inflammatory syndrome (rare, following COVID-19)
\end{itemize}

\section*{Prognosis and Outcomes}
Most people with common coronavirus infections recover within one to two weeks. For COVID-19, the vast majority of infections are mild or moderate, but older people and those with chronic conditions are at higher risk of severe illness, hospitalisation, and death. Vaccination and antiviral treatment have markedly improved outcomes. A minority of people develop long COVID with symptoms persisting for months. MERS and SARS carry higher fatality rates, but these viruses currently cause only limited, sporadic disease.

\section*{Prevention and Self-Care}
\begin{itemize}
    \item Stay up to date with COVID-19 and influenza vaccinations
    \item Wash hands regularly and use hand sanitiser
    \item Cover coughs and sneezes, and dispose of tissues safely
    \item Stay home when unwell
    \item Consider masks in crowded indoor settings during outbreaks, especially if at high risk
    \item Improve indoor ventilation by opening windows
    \item Seek early assessment for antiviral treatment if you are in a high-risk group
    \item Maintain a healthy lifestyle to support the immune system
\end{itemize}

\section*{Sources and Further Reading}
The following reputable sources provide further information for healthcare students and patients seeking reliable, current advice about coronaviruses and COVID-19.

\begin{itemize}
    \item World Health Organization. \textit{Coronavirus disease (COVID-19) fact sheets.} https://www.who.int/
    \item Australian Government Department of Health and Aged Care. \textit{COVID-19 information.} https://www.health.gov.au/
    \item Healthdirect Australia. \textit{COVID-19.} https://www.healthdirect.gov.au/covid-19
    \item Centers for Disease Control and Prevention. \textit{Coronavirus (COVID-19).} https://www.cdc.gov/
    \item NHS. \textit{Coronavirus (COVID-19).} https://www.nhs.uk/conditions/coronavirus-covid-19/
    \item MedlinePlus. \textit{Coronavirus infections.} https://medlineplus.gov/coronavirusinfections.html
\end{itemize}
